\documentclass[11pt]{article}
\usepackage[margin=0.92in]{geometry}
\usepackage{amsmath,amssymb,amsthm,mathtools}
\usepackage{graphicx,microtype,xcolor}
\usepackage[colorlinks=true,linkcolor=blue!55!black,citecolor=blue!55!black,urlcolor=blue!55!black]{hyperref}
\emergencystretch=1em

\newtheorem{theorem}{Theorem}
\newtheorem{corollary}{Corollary}
\newcommand{\R}{\mathbb R}
\newcommand{\N}{\mathbb N}
\newcommand{\comp}{\operatorname{comp}}
\newcommand{\diam}{\operatorname{diam}}

\title{Global Bang Degree Does Not Control Sublevel Topology:\\A Level-Centered Repair}
\author{Automated mathematics run \texttt{fa\_banach\_001}, lane 7}
\date{17 August 2026}

\begin{document}
\maketitle

\begin{center}
\fbox{\parbox{0.92\textwidth}{\textbf{Status.} Substantial partial result and
sharp obstruction, likely valid and pending human review.  A fixed global Bang
degree cannot bound the affine-section component constants of quasianalytic
sublevel sets, even in a fixed analytic class on a fixed rectangle.  Bang
degrees of the shifted level functions $f-t$ and $f+t$ do give an explicit
line-section bound.}}
\end{center}

\section{Source prompt}

For a sublevel set
\[
 K_t=\{x\in K:|f(x)|\le t\},
\]
Rainer proves a metric-entropy estimate containing a constant
\[
 C_{K_t}=C(n)\max_{1\le j\le n-1}\omega_j B_{n-j}(K_t),
\]
where $B_{n-j}(K_t)$ bounds the number of connected components of
$P\cap K_t$ over affine $(n-j)$-planes $P$.  Definability gives finiteness, and
the source asks for explicit bounds ``via the Bang degree depending on the
Denjoy--Carleman complexity of $K_t$''
\cite[after Corollary~1.18]{Rainer}.

\begin{figure}[ht]
\centering
\includegraphics[width=0.94\textwidth]{figures/open_problem_crop.png}
\caption{The metric-entropy corollary and the prompt on source PDF page 9.}
\end{figure}

The source does not define a numerical Denjoy--Carleman complexity or format
for a definable presentation of $K_t$.  The results below isolate a necessary
scale that any such complexity must retain.

\section{Proof intuition}

Derivative bounds permit very high frequency when its amplitude is
exponentially small.  Adding that oscillation to a nonzero constant leaves the
global size of the function, and therefore its global Bang degree, essentially
unchanged.  But taking a sublevel exactly at the constant term removes the
large offset and exposes every oscillation topologically.

This also identifies the repair.  On an interval, the boundary of
$\{|f|\le t\}$ consists of zeros of the two shifted functions $f-t$ and
$f+t$.  The source's quasianalytic zero theorem applies to those shifts, so the
relevant degrees must be centered at the boundary levels.  Their norms detect
the tiny oscillatory amplitude that the global degree misses.

\section{Notation from the source}

For an admissible weight $\mu=(\mu_j)_{j\ge1}$, base derivative bound $M>0$,
length scale $a>0$, and lower norm bound $b>0$, write
\[
 \mathcal N(\mu,M,a,b)
 :=2D_{a\mu}(b/M)
\]
for the associated degree of the source, whenever the quadruple is admissible.
For quasianalytic $\mu$ it is always finite.  Proposition~10.1 of the source
states that a $(\mu,M)$-smooth function on an interval of length $a$ whose
supremum norm is at least $b$ has at most
$\mathcal N(\mu,M,a,b)$ zeros, counted with multiplicity
\cite[Proposition~10.1]{Rainer}.

\section{A fixed-class obstruction}

\begin{theorem}[global degree cannot control topology]\label{thm:negative}
Let
\[
 K=[0,2\pi]\times[0,1],\qquad \mu_j=j,\qquad M_0=2.
\]
For every integer $N\ge1$, there is a restricted analytic
$(\mu,M_0)$-smooth function $f_N:K\to\R$ such that the associated global
degree with lower bound $b_0=1$ is independent of $N$, while
\[
 B_1(\{x\in K:|f_N(x)|\le1\})\ge N+1.
\]
In particular, no bound depending only on the global Bang degree, the fixed
weight and derivative majorant, the dimension, and the fixed geometry of $K$
can control the component constants in the source prompt.
\end{theorem}

\begin{proof}
Set
\[
                  f_N(x,y)=1+e^{-N}\sin(Nx).
\]
For $j\ge1$, only the pure $x$-derivative contributes to the source's
multivariate derivative seminorm, and
\begin{equation}\label{eq:derivative}
                  \|f_N\|_{j,K}=e^{-N}N^j\le j!\le2j!.
\end{equation}
Indeed, $j!\ge(j/e)^j$, and with $u=j/N$ the desired first inequality follows
from
\[
 \left(\frac{j}{eN}\right)^j\ge e^{-N}
 \quad\Longleftrightarrow\quad
 N(u\log u-u+1)\ge0.
\]
Also $\|f_N\|_K=1+e^{-N}<2$.  Since the weight $\mu_j=j$ gives
$M_j=2j!$, equation \eqref{eq:derivative} proves
$(\mu,2)$-smoothness for every $N$.  The fixed lower bound
$\|f_N\|_K\ge1$ therefore uses the same associated degree
\[
                  \mathcal N(\mu,2,\diam K,1)
\]
for the whole family.  The degree obtained by inserting the actual norm
$\|f_N\|_K$ is no larger, by monotonicity in the lower bound.

Because $1-e^{-N}>0$, the level-one sublevel condition is
\[
 |f_N(x,y)|\le1\quad\Longleftrightarrow\quad\sin(Nx)\le0.
\]
On $[0,2\pi]$ this set is the isolated point $\{0\}$ together with the $N$
disjoint intervals
\[
 \left[\frac{(2k+1)\pi}{N},\frac{(2k+2)\pi}{N}\right],
 \qquad k=0,\ldots,N-1.
\]
Their products with $[0,1]$ are $N+1$ connected components.  Intersecting
with any horizontal line gives the same $N+1$ components, proving the asserted
lower bound for $B_1$.
\end{proof}

\section{The level-centered repair}

\begin{theorem}[explicit line-section bound]\label{thm:positive}
Let $I\subset\R$ be a nondegenerate compact interval, let $\mu$ be a
quasianalytic admissible weight, and let $g:I\to\R$ be
$(\mu,M_0)$-smooth.  Fix $t>0$ and put
\[
 \widehat M_0=2(M_0+t),\qquad
 b_\sigma=\|g-\sigma t\|_I\quad(\sigma\in\{-1,1\}).
\]
If $g$ is not identically $t$ or $-t$, then
\begin{equation}\label{eq:component-bound}
 \comp\{s\in I:|g(s)|\le t\}
 \le1+\sum_{\sigma\in\{-1,1\}}
 \mathcal N(\mu,\widehat M_0,|I|,b_\sigma).
\end{equation}
If $g\equiv t$ or $g\equiv-t$, the sublevel set is $I$ and has one component.
\end{theorem}

\begin{proof}
Assume $g$ is not identically either boundary level.  The boundary in $I$ of
$S=\{|g|\le t\}$ is contained in
\[
                  \{g-t=0\}\cup\{g+t=0\}.
\]
Consequently, if $Z_I(h)$ denotes the number of distinct zeros of $h$ in $I$,
\begin{equation}\label{eq:topological}
                       \comp(S)\le Z_I(g-t)+Z_I(g+t)+1.
\end{equation}
Both shifted functions are $(\mu,\widehat M_0)$-smooth: their zeroth norms
are at most $M_0+t<\widehat M_0$, and all positive-order derivatives are
those of $g$.  Moreover $b_\sigma>0$.  The source's zero theorem bounds
$Z_I(g-\sigma t)$ by
$\mathcal N(\mu,\widehat M_0,|I|,b_\sigma)$.  Substitution in
\eqref{eq:topological} proves \eqref{eq:component-bound}.
\end{proof}

\begin{corollary}[uniform planar entropy constant]\label{cor:plane}
Let $K\subset\R^2$ be a convex body and let $f$ be
$(\mu,M_0)$-smooth on a convex neighborhood of $K$.  Fix $t>0$.  Suppose
there is $b_*>0$ such that for every affine line $P$, every nondegenerate
interval $I=P\cap K$, and both $\sigma\in\{-1,1\}$, either
$f|_I\equiv\sigma t$ or
\[
                         \|f|_I-\sigma t\|_I\ge b_*.
\]
Then, with
\[
 \mathcal N_*=
 \mathcal N(\mu,2(M_0+t),\diam K,b_*),
\]
one has
\[
                 B_1(K_t)\le2\mathcal N_*+1.
\]
Thus the component term in the constant of Corollary~1.18 of the source may
be taken at most $C(2)\omega_1(2\mathcal N_*+1)$.
\end{corollary}

\begin{proof}
Every nonempty line section of a convex body is an interval.  Parameterizing
it by arclength preserves the derivative majorant.  Apply
Theorem~\ref{thm:positive}; use $|I|\le\diam K$, the assumed lower bound, and
monotonicity of $\mathcal N$ in its length and lower-bound arguments.  A
restriction identically equal to a boundary level contributes one component
and obeys the same final bound.
\end{proof}

\section{Scope, verification, and novelty audit}

The construction and estimates are exact and have no computational
dependency.  The elementary inequality in \eqref{eq:derivative} is uniform in
both $j$ and $N$.  The isolated endpoint component at $x=0$ is included, so
the count is $N+1$, not merely $N$.  In the same family,
$\|f_N-1\|_K=e^{-N}$; hence the level-centered associated degree is at least
twice its starting index, which grows linearly with $N$.  It therefore records
exactly the scale erased by the constant offset.

This does not solve the source's full higher-dimensional request.  The phrase
``Denjoy--Carleman complexity'' is not assigned a format or numerical value in
the paper, and scalar line restrictions do not bound the number of components
of higher-dimensional affine sections.  Such a result would require effective
control of multivariate critical systems or an effective stratification
theorem for a specified Denjoy--Carleman presentation.

A bounded novelty audit checked the run's result, attempt, and proof-gap
indexes, arXiv:2606.08183, the exact metric-entropy prompt, and the zero-count
result cited in the source.  No prior packet or later resolution was found.
The source is very recent and a final external-source refresh was unavailable,
so novelty remains pending specialist review.

Recommended human-review focus is: (1) whether the intended notion of
Denjoy--Carleman format was meant to include a separate frequency-sensitive
presentation parameter; (2) the fixed-weight derivative estimate; and (3)
the passage from the two shifted zero counts to the planar uniform constant.

\begin{thebibliography}{9}
\bibitem{Rainer}
A.~Rainer,
\emph{Effective quasianalytic Remez inequalities on tame sets},
arXiv:2606.08183 (2026).
\end{thebibliography}

\end{document}
